\section{The rank-one matrix mechanism}\label{mat:sec:matrix}
In this section matrices are real and finite dimensional. For $A\succ0$,
$B=A+uu^T$, and $0<\alpha<1$, define
\[
 Q_\alpha(A,u)=A^{-\alpha/2}B^\alpha A^{-\alpha/2},\qquad
 C_\alpha(A,u)=Q_\alpha(A,u)-I.
\]
The leverage is $\tau=u^TA^{-1}u$. Powers are principal powers.
The eigenvalues of $Q_\alpha$ are the generalized eigenvalues of
$(B^\alpha,A^\alpha)$. This is an energy-relative comparison; it is
different from dividing $\norm{B^\alpha-A^\alpha}$ by its own norm.

For connected graph Laplacians, all inverse powers in this paper mean
the Moore--Penrose functional calculus: zero on the common kernel and
the stated negative power on its orthogonal complement. The relative
matrix $C$ is zero on that kernel. Every determinant on the common
support is explicitly understood as such.

\begin{theorem}[Complete closure of rank-one relative spectra]
\label{mat:thm:closure}
Fix an integer $n\ge1$, $0<\alpha<1$, and $\tau\ge0$.
The closure of the ordered eigenvalue lists of
$Q_\alpha(A,u)$ over all $A\in\R^{n\times n}$ positive definite and
all $u\in\R^n$ with $u^TA^{-1}u=\tau$ is exactly
\[
 \mathcal S_{n,\alpha,\tau}
 =\left\{(q_1,\ldots,q_n):
       q_1\ge\cdots\ge q_n\ge1,\quad
       \prod_{i=1}^nq_i=(1+\tau)^\alpha\right\}.
\]
For $\tau=0$ this set is the singleton $(1,\ldots,1)$.
The closure is essential: the theorem asserts approximation of each
list, without asserting that every list is attained at finite condition
number.
\end{theorem}

\begin{corollary}[Exact worst-case strict count]\label{mat:cor:count}
For $\tau>0$ and $\varepsilon>0$, put
\[
 R=\frac{\alpha\log(1+\tau)}{\log(1+\varepsilon)},\qquad
 M=\max\{0,\lceil R\rceil-1\}.
\]
Then
\[
 \sup_{A\succ0,\ u^TA^{-1}u=\tau}
 \#\{j:\lambda_j(C_\alpha(A,u))>\varepsilon\}=M,
\]
where the supremum ranges over all finite dimensions.
In a fixed dimension $n$, the exact supremum is $\min(n,M)$.
For $\tau=0$ both suprema are zero.
\end{corollary}

For fixed $\tau>0$, the worst-case count is
\[
 M=\frac{\alpha\log(1+\tau)}{\varepsilon}+O_{\alpha,\tau}(1)
 \quad(\varepsilon\downarrow0).
\]
Thus no uniform logarithmic-in-accuracy rank bound is possible in this
relative metric, even for a single positive rank-one update of fixed
leverage. This does not contradict absolute-error compression estimates.

\subsection{Separated scales and the spectral closure}
The following integral supplies the whole construction.
Write $c_\alpha=\sin(\pi\alpha)/\pi>0$.

\begin{lemma}[Separated-scale rank-one limit]\label{mat:lem:scales}
Let $r_1,\ldots,r_n\ge1$, and set
\[
 z_i=(r_i-1)\prod_{j=i+1}^n r_j,\qquad
 a_i=\delta^{n-i},\quad
 A_\delta=\diag(a_1,\ldots,a_n),\quad
 (u_\delta)_i=\sqrt{a_i z_i},\qquad 0<\delta<1.
\]
Then, for every $\delta$,
\[
 u_\delta^TA_\delta^{-1}u_\delta
    =\sum_i z_i=\prod_i r_i-1.
\]
Moreover, in operator norm as $\delta\downarrow0$,
\[
 C_\alpha(A_\delta,u_\delta)
       \longrightarrow\diag(r_1^\alpha-1,\ldots,r_n^\alpha-1).
\]
\end{lemma}
\begin{proof}
The telescoping identity for $\sum_i z_i$ is immediate from
$z_i=\prod_{j=i}^n r_j-\prod_{j=i+1}^n r_j$.
The resolvent formula for fractional powers and the rank-one inverse
identity give
\begin{equation}\label{mat:eq:resolvent}
 B^\alpha-A^\alpha
 =c_\alpha\int_0^\infty
 \frac{t^\alpha(A+tI)^{-1}uu^T(A+tI)^{-1}}
 {1+u^T(A+tI)^{-1}u}\,dt.
\end{equation}
The integral converges in norm; its integrand is positive semidefinite.
For its $i$th relative diagonal entry, substitute $t=a_i s$ to obtain
\begin{equation}\label{mat:eq:diag}
 C_{ii}=c_\alpha z_i\int_0^\infty
 \frac{s^\alpha}{(1+s)^2}
 \left(1+\sum_j\frac{a_jz_j}{a_j+a_i s}\right)^{-1}ds.
\end{equation}
For every $s>0$, $a_j/a_i\to0$ if $j<i$, and $a_j/a_i\to\infty$
if $j>i$. The denominator in parentheses therefore tends to
\[
 H_i+\frac{z_i}{1+s},\qquad
 H_i=1+\sum_{j>i}z_j=\prod_{j>i}r_j.
\]
The integrable function $s^\alpha/(1+s)^2$ dominates the integrand
after removing the fixed multiplier $c_\alpha z_i$. Dominated
convergence gives
\[
 C_{ii}\longrightarrow
 c_\alpha\frac{z_i}{H_i}
    \int_0^\infty\frac{s^\alpha}{(1+s)(1+s+z_i/H_i)}\,ds
 = (1+z_i/H_i)^\alpha-1=r_i^\alpha-1.
\]
The scalar equality follows by applying the same fractional resolvent
formula to the scalars $1+z_i/H_i$ and $1$. It also holds when $z_i=0$.

For $i<j$, omit the positive denominator in \eqref{mat:eq:resolvent}.
The resulting scalar divided-difference integral yields
\begin{align*}
 |C_{ij}|
 &\le\sqrt{z_i z_j}(a_i a_j)^{(1-\alpha)/2}
                  \frac{a_j^\alpha-a_i^\alpha}{a_j-a_i}\\
 &=\sqrt{z_i z_j}\,x^{(1-\alpha)/2}
                    \frac{1-x^\alpha}{1-x},\qquad x=a_i/a_j.
\end{align*}
This tends to zero because $0<\alpha<1$ and $x\to0$.
Symmetry treats $i>j$. Entrywise convergence in fixed dimension implies
operator-norm convergence.
\end{proof}

\begin{proof}[Proof of Theorem~\ref{mat:thm:closure}]
Formula \eqref{mat:eq:resolvent} proves $B^\alpha-A^\alpha\succeq0$,
so every eigenvalue of $Q_\alpha$ is at least one. The determinant
lemma gives
\[
 \det Q_\alpha
    =\left(\frac{\det(A+uu^T)}{\det A}\right)^\alpha
    =(1+\tau)^\alpha.
\]
These constraints are closed. Conversely, given an admissible list,
take $r_i=q_i^{1/\alpha}$ and apply Lemma~\ref{mat:lem:scales}.
Its leverage is exactly $\tau$ at every $\delta$, and its relative
eigenvalues converge to the desired list, after ordering.
For $\tau=0$, positive definiteness forces $u=0$.
\end{proof}

\begin{proof}[Proof of Corollary~\ref{mat:cor:count}]
If $m$ correction eigenvalues exceed $\varepsilon$, then
\[
 (1+\varepsilon)^m<\prod_j(1+\lambda_j(C))
                 =(1+\tau)^\alpha,
\]
so $m<R$. The largest allowable integer is $M$; in dimension $n$ it
is also at most $n$.
For the reverse inequality, let $m=\min(n,M)>0$. Choose a target
list with $m$ entries equal to $(1+\tau)^{\alpha/m}>1+\varepsilon$
and the remaining $n-m$ entries equal to one. Theorem~\ref{mat:thm:closure}
approximates it closely enough to retain those $m$ strict inequalities.
For the unrestricted supremum take dimension $M$ when $M>0$.
The case $M=0$ is immediate. At integral $R$, the strict inequality
excludes $m=R$, which explains $\lceil R\rceil-1$ rather than
$\lfloor R\rfloor$.
\end{proof}

